\begin{figure}[htbp]
\centering
\begin{tikzpicture}[
  font=\fontsize{9}{11}\selectfont,
  % Заглавие на жизнена линия.
  head/.style={box, fill=gray!10, minimum height=0.8cm},
  % Жизнените линии са пунктирани, за да не се четат като поток от данни.
  life/.style={densely dotted, gray!60},
  msg/.style={-Stealth, thick},
  ret/.style={-Stealth, semithick, densely dashed},
  lbl/.style={midway, above, font=\fontsize{8}{9.6}\selectfont, align=center, inner sep=1.5pt},
  lblb/.style={midway, below, font=\fontsize{8}{9.6}\selectfont, align=center, inner sep=1.5pt},
  tiny note/.style={font=\fontsize{7.5}{9}\selectfont, align=left, inner sep=1.5pt},
]

% === жизнени линии ===========================================================
% Третата линия носи две имена, защото read_prefix чете, а после връща обвивката
% PrefaceConnection. По потока това е един и същ участник.
\node[head, minimum width=2.1cm] (hA) at (0,0)    {Клиент};
\node[head, minimum width=2.1cm] (hB) at (3.0,0)  {Сокет};
\node[head] (hC) at (7.2,0)  {\code{read\_prefix}\\ \code{PrefaceConnection}};
\node[head] (hD) at (11.2,0) {Протоколен слой\\ (TLS ръкостискане\\ или HTTP парсер)};

\draw[life] (0,-0.62)    -- (0,-8.80);
\draw[life] (3.0,-0.62)  -- (3.0,-8.80);
\draw[life] (7.2,-0.62)  -- (7.2,-8.80);
\draw[life] (11.2,-0.92) -- (11.2,-8.80);

% === четене и решение ========================================================
\draw[msg] (0,-1.55) -- (3.0,-1.55)
  node[lbl] {изпраща \code{0x16 ...}\\ или \code{GET ...}};

\draw[msg] (7.2,-2.35) -- (3.0,-2.35) node[lbl] {\code{async\_read}};
\draw[ret] (3.0,-3.05) -- (7.2,-3.05) node[lbl] {първите октети};

% Възловият момент. Октетите вече са извън ядрото, затова обвивката трябва да ги
% пази. Кутията стои между линиите на Сокет и Протоколен слой, за да не ги закрие.
\node[box, fill=gray!25, align=center, text width=6.1cm,
      font=\fontsize{8}{10}\selectfont] (wrap) at (7.2,-4.05)
  {\code{classify} решава по първия октет\\
   сокетът се обвива в \code{PrefaceConnection},\\
   която пази прочетените октети};

\draw[msg] (7.2,-5.35) -- (11.2,-5.35) node[lbl] {предава връзката нататък};

% === повторение на прочетеното ===============================================
% Щрихованата лента показва живота на буфера: започва при обвиването и свършва
% при първото четене на протоколния слой. Стои вляво от жизнената линия, за да
% не закрива върховете на стрелките, които идват отдясно.
\draw[pattern=north east lines] (6.83,-4.85) rectangle (7.09,-6.80);
\draw[semithick] (6.83,-5.70) -- (5.6,-5.70);
\node[tiny note, anchor=east, align=right] at (5.55,-5.70)
  {буфер с\\ прочетените\\ октети};

\draw[msg] (11.2,-5.95) -- (7.2,-5.95) node[lbl] {\code{async\_read}};
\draw[ret] (7.2,-6.85) -- (11.2,-6.85)
  node[lbl] {повторените октети,\\ без системно извикване};

% === продължението от сокета =================================================
\draw[msg] (11.2,-7.45) -- (7.2,-7.45) node[lbl] {\code{async\_read}};
\draw[msg] (7.2,-7.45) -- (3.0,-7.45) node[lbl] {препраща}
                                      node[lblb] {буферът е празен};

\draw[ret] (3.0,-8.35) -- (7.2,-8.35);
\draw[ret] (7.2,-8.35) -- (11.2,-8.35) node[lbl] {останалото от потока};

% === какво вижда протоколният слой ===========================================
% Осемте клетки са схема на потока. Границата между щрихованите и сивите клетки
% умишлено не е подчертана, защото протоколният слой не я вижда.
\foreach \i in {0,1,2} {
  \draw[pattern=north east lines] (0.4+0.52*\i,-10.10) rectangle (0.92+0.52*\i,-9.60);
}
\foreach \i in {3,4,5,6,7} {
  \draw[fill=gray!10] (0.4+0.52*\i,-10.10) rectangle (0.92+0.52*\i,-9.60);
}

\draw[decorate, decoration={brace, amplitude=4pt}] (0.4,-9.50) -- (1.96,-9.50)
  node[midway, above, font=\fontsize{7.5}{9}\selectfont, inner sep=5pt] {повторени};
\draw[decorate, decoration={brace, amplitude=4pt}] (1.96,-9.50) -- (4.56,-9.50)
  node[midway, above, font=\fontsize{7.5}{9}\selectfont, inner sep=5pt] {от сокета};
\draw[decorate, decoration={brace, amplitude=4pt, mirror}] (0.4,-10.20) -- (4.56,-10.20)
  node[midway, below, align=center, font=\fontsize{8}{10}\selectfont, inner sep=5pt]
  {един непрекъснат поток\\ за протоколния слой};

% Шрифтът е записан тук, а не чрез стила note, защото преамбюлът предефинира
% note по-нататък и кутията става твърде голяма, а тогава закрива стрелката
% с останалото от потока.
\node[grp, align=left, font=\fontsize{8}{10}\selectfont,
      text width=6.6cm, anchor=west] at (5.70,-9.90)
  {Защо не \code{MSG\_PEEK}: октетите, прочетени за разпознаването,
   са същите, които следващият слой ще поиска. Надничането ги чете
   два пъти и иска отделна реализация във всеки бекенд.};

\end{tikzpicture}
\caption{Четене с връщане на прочетеното на пътя на приемане.
\code{read\_prefix} взема първите октети от сокета и \code{classify} решава по
първия от тях. Сокетът се обвива в \code{PrefaceConnection}, която пази
прочетеното. Протоколният слой получава първо повторените октети, после
останалото от сокета, и вижда потока така, все едно никой не го е поглеждал.}
\label{fig:pushback_connection}
\end{figure}
